\section{Sharp numerical bounds from the signature}\label{sec:numerical}

A numerical component of height \(h\) and dimension \(m\) is
a convex polygon with slopes in \([0,1]\), endpoints \((0,0)\)
and \((h,m)\), and integral breakpoints. Equivalently, each
reduced slope \(a/c\) occurs with multiplicity divisible by \(c\)
\cite[\S1.1]{oort}. The polygon order used here is pointwise order
on the prefix sums \(S_j\) of the ordered slopes. Duality sends
\(S_j\) to \(j-m+S_{h-j}\), so it preserves this order.

Orient the Hodge multiplicities as \((m_i,h-m_i)\) with
\(m_i\le h/2\), and let \(b\) be the number of positive \(m_i\).
The ordinary slack at the \(j\)-th norm bound, before the strict
unit margin, is
\begin{equation}\label{eq:ordinary-slack}
 D_j=\sum_i\min(j,h-j,m_i).
\end{equation}
Indeed the ordinary two blocks have a total rounded prefix of
\(j-\min(j,h-j,m_i)\). If \(b>0\), an excess \(\Delta_j\) above
the ordinary rounded prefix is allowed precisely when
\(\Delta_j\le D_j-1\). The all-definite case is ordinary immediately.

\begin{proof}[Proof of Corollary~\ref{cor:ranktwo}]
The sum of the two slopes on each balanced side is one.
An ordinary balanced pair has zero excess at the only index
\(j=1\). For a nonordinary pair both slopes on both sides are
strictly positive, so each rounded first prefix is one and the
excess is two. As the ordinary slack is \(b\), the strict margin
gives \(2r\le b-1\). Each such pair removes two units of
\(p\)-rank, giving the stated bound. With no balanced pair,
each split component is definite and ordinary.
\end{proof}

\begin{proof}[Proof of Corollary~\ref{cor:rankfour}]
In a nonordinary pair of type \((2,2)\), each side has fewer than
two zero slopes, since two zero slopes would force the other
two to be one. The sum of the two rounded prefixes at \(j=2\)
is therefore at least two on this pair, compared with zero in
the ordinary case. Its slack is two, and definite pairs add
no slack. This violates the strict unit margin.
\end{proof}

\begin{theorem}[Least polygon with fixed endpoint multiplicities]\label{thm:extremal}
Fix \(h,m\), and fix the numbers \(z,o\) of zero and unit slopes
of a numerical component. Set \(q=h-z-o\), \(w=m-o\).
If \(q=0\) the polygon is ordinary. Otherwise \(0<w<q\).
The least possible interior polygon is the chord from \((0,0)\) to
\((q,w)\) when \(w=1\) or \(q-w=1\). In the remaining cases it
joins \((0,0)\), \((q-w,1)\), and \((q,w)\), with interior slopes
\[
 \left(\frac{1}{q-w}\right)^{q-w},\qquad
 \left(\frac{w-1}{w}\right)^w .
\]
Consequently the least possible paired excess at each \(j\)
is attained simultaneously and is
\begin{equation}\label{eq:interval}
  \mathbf 1_{z<j<h-o}+\mathbf 1_{o<j<h-z}.
\end{equation}
The same expression is zero for ordinary endpoints.
\end{theorem}

\begin{proof}
Write \(a=q-w\) and \(c=w\). Each proper interior vertex
\((x,y)\) has integer coordinates with \(y\ge1\) and
\(y\ge x-a+1\). The second bound follows because the remaining
slopes are strictly less than one. If \(c=1\) there is no such
vertex because its ordinate would be strictly between zero and one.
If \(a=1\) the same holds after dualizing.

For \(a,c\ge2\), the displayed two slopes are increasing, with
equality possible when \(a=c=2\), in which case the middle point is
collinear. The resulting lower convex polygon is the maximum of
the lines \(x/a\) and \(x-a+1-(x-a)/c\).
At any proper vertex the first line is bounded by \(1\) if
\(x\le a\), and by \(x-a+1\) if \(x\ge a\).
The second line is bounded by \(1\) when \(x\le a\), and by
\(x-a+1\) otherwise. Thus both lie below every vertex. Convexity
puts the proposed polygon below every joining segment, including
the isoclinic chord.

Restore the zero and unit tails. From \(j=z+1\) through
\(j=h-m=z+a\), the first interior segment has rounded prefix one.
At \(j=h-m+t\), \(1\le t<c\), the next segment has rounded prefix
\(1+t\), while at \(t=c\) the prefix equals \(c\). Hence excess
over the ordinary prefix is one precisely for \(z<j<h-o\).
If \(a=1\) or \(c=1\), taking the ceiling of the chord gives
the same strict interval. Dualizing swaps \(z,o\). Since duality
preserves order, the minimizer supplies both terms in
\eqref{eq:interval} simultaneously.
\end{proof}

The symmetric specialization of this polygon, with equal numbers
of extreme slopes and equal residual dimension and codimension,
appears in Achter and Pries \cite[Definition~2.4]{achterpries}.
The oriented component statement gives an endpoint-only feasibility
test for the strict simultaneous bounds. For each pair select
either the ordinary endpoints or integers
\(0\le z\le h-m-1\), \(0\le o\le m-1\).
Test the sum of \eqref{eq:interval} against \(D_j-1\) and minimize
\(\sum_i(z_i+o_i)\). Any feasible endpoints lift to numerical
polygons by Theorem~\ref{thm:extremal}, and every admissible
collection gives feasible endpoints. This proves exactness.

\begin{proposition}[When the inequalities force ordinary reduction]\label{prop:signatures}
The only signatures for which the numerical restrictions force the
ordinary polygon are all-definite signatures, one mixed pair of type
\((1,h-1)\), two mixed pairs in rank two, and one mixed pair of
type \((2,2)\) in rank four. Definite pairs can be appended freely
to all cases.
\end{proposition}

\begin{proof}
If a lone mixed pair has \(m=1\), then at \(j=1\) the slack
is one. A nonordinary block on its dimension-one side cannot
have a unit slope, so the complementary block has no zero, and
the paired rounded prefix increases by at least one.
The rank-two and rank-four exceptions follow from the arguments
for Corollaries~\ref{cor:ranktwo} and~\ref{cor:rankfour}.

For the converse, in one mixed pair replace the last zero and
first unit slope by two copies of \(1/2\), preserving duality.
This is the least nonordinary paired polygon for that pair.
In fact a nonordinary block with endpoint losses \(a,c\ge1\) has
prefix at the ordinary kink \(h-m\) at least \(1/2\) by
Theorem~\ref{thm:extremal}, while its other integer prefixes
are at least ordinary. The replacement has excess
\[
        \Delta_j=\mathbf 1_{j=m}+\mathbf 1_{j=h-m}
\]
by \eqref{eq:interval}. For \(h=2\) and \(b\ge3\), the slack is
at least three at its only index, so the replacement is allowed.
For \(h>2\) and at least two mixed pairs, choose any mixed pair.
If its two affected indices are distinct, the slack is at least
two at each. If they coincide, then \(m=h/2\ge2\), and at that
index the modified pair contributes at least two to the slack
while a second pair contributes at least one. Finally if there is
one mixed pair with \(m\ge2\), the slack at both affected indices
is \(m\). The replacement fails the strict unit margin precisely
when they coincide at \(m=2\), the already listed rank-four case.
This gives a witness in every other signature.
\end{proof}

The total number of signatures up to dimension \(G\) is
\[
 \sum_{h=1}^G
 \left(\binom{\lfloor h/2\rfloor+\lfloor G/h\rfloor+1}
                    {\lfloor h/2\rfloor+1}-1\right)
\]
by counting multisets of \(e\) multiplicities and summing over \(e\).
For \(G=30\) this is 1188. The four disjoint cases of
Proposition~\ref{prop:signatures} give 111, 81, 14 and 7,
respectively. The first count is
\(\sum_{h=1}^{30}\lfloor30/h\rfloor\), the second omits \(h=1\),
and the remaining cases count \(2\le e\le15\) and \(1\le e\le7\).

\begin{theorem}[Sharp numerical \(p\)-rank]\label{thm:sharp-prank}
If \(b=0\), the minimum is \(eh\). Otherwise put
\[
 H=\lfloor h/2\rfloor,\qquad L_j=\max(0,2b+1-D_j)
 \quad (1\le j\le H).
\]
The sharp minimum is \(h(e-b)+M\), where the mixed contribution
\(M\) is as follows.
\begin{itemize}
\item If all mixed multiplicities are one, then
\[
 M=\begin{cases}
  2\lceil(b+1)/2\rceil, & h=2,\\
  h+H(b-1), & h\text{ odd},\\
  h+(H-1)(b-1)+\lceil(b-1)/2\rceil, & h\ge4\text{ even}.
 \end{cases}
\]
\item If \(h=4\) and all mixed multiplicities are two, then \(M=b+3\).
\item In every other case with some mixed multiplicity at least two,
\[
 M=\begin{cases}
  \displaystyle\sum_{j=1}^H L_j, & h\text{ odd},\\
  \displaystyle\sum_{j=1}^{H-1}L_j+\lceil L_H/2\rceil,
      & h\text{ even}.
 \end{cases}
\]
\end{itemize}
\end{theorem}

\begin{proof}
For a mixed pair write \(P_j=2-\Delta_j\), using the minimizing
endpoints. The simultaneous condition is
\(\sum P_j\ge L_j\) for \(1\le j\le H\).
For a nonordinary pair, \(o<H\) and
\[
 P_j=\mathbf 1_{j\le z}+\mathbf 1_{j\le o}
             +\mathbf 1_{j\ge h-z}.
\]
For an ordinary pair, \(P_j=2\).
If \(h\) is odd, summing \(P_j\) gives exactly \(z+o\) for a
nonordinary pair and \(h-1\) for an ordinary pair. If \(h\) is even,
give the index \(H\) weight one half. This weighted sum has the
same two values, and \(P_H\) is always even.
These observations give the lower bounds in the last case.

For sharpness there, use only nonordinary pairs. Think of the two
endpoint counts as slots of capacities
\(\min(H,h-m_i-1)\) and \(m_i-1\).
Choose lengths in the slots with prescribed tail counts \(L_j\)
for \(j<H\), and \(L_H\) in odd rank or \(\lceil L_H/2\rceil\)
in even rank at \(H\). At these lengths the third indicator in
\(P_j\) is absent except at the even midpoint, where it doubles
the contribution of a slot of length \(H\).
For a decreasing list of tail counts a feasible assignment
exists if each count is at most the number of slots with capacity
at least its index. To see this, sort the lengths and capacities
and match in decreasing order.

At \(j=1\) there are \(b\) zero slots and at least one unit slot.
For \(2\le j<H\), and for \(j=H\) in odd rank, the \(b\)
zero slots suffice since \(D_j\ge b+1\).
At the even midpoint only unbalanced pairs supply a zero slot.
Let their number be \(u\). For \(H\ge3\), if there is any
balanced pair, then \(D_H\ge u+3(b-u)\) and \(L_H\le u\).
If there is none, \(u=b\) already suffices. For \(H=2\) in
this case \(u\ge1\), \(L_H=u+1\), and
\(\lceil(u+1)/2\rceil\le u\). The resulting assignment
attains the lower bounds.

If all mixed multiplicities are one and \(h\ge3\), at least
one pair is ordinary. Otherwise all \(o=0\) and each \(P_1\)
is at most one, whereas \(L_1=b+1\).
With \(c\ge1\) ordinary pairs, each index needs at least
\(T=\max(0,b+1-2c)\) further units of \(P_j\).
In odd rank choose \(T\) remaining zero endpoints of length \(H\).
In even rank use \(\lceil T/2\rceil\) of length \(H\) and the
remaining \(T-\lceil T/2\rceil\) of length \(H-1\).
There are enough slots because \(T\le b-c\), and they attain
the weighted lower bound for this \(c\). Increasing \(c\) costs
\(h\) and can save at most \(h-1\) of this lower bound, so
\(c=1\) is optimal. Rank two follows directly by counting
ordinary pairs.

Finally in rank four with all mixed multiplicities two,
every nonordinary pair has \(P_2=0\), so at least one
ordinary pair is necessary. With \(c\) such pairs,
the lower bound is \(4c+\max(0,b+1-2c)\), from index one.
It is minimized at \(c=1\), and is attained with one endpoint
of length one in each of the other pairs.
\end{proof}

The dimension-30 table includes exact polygon multiplicities
as well as endpoint minima. The endpoint formula gives feasibility
and optimization. Individual polygons with fixed endpoints can
differ in admissibility. For instance, at \(h=8,m=4,z=2,o=1\)
the extremal interior \((1/2)^2,(2/3)^3\) passes with one mixed
pair, while the alternative \((3/5)^5\) fails at \(j=4\).
For full multiplicities one can transfer the paired rounded
profiles of integral convex paths instead. The verification
replays retain both the numerical minima and the historical full
counts. The closed formula gives minima for \(5{,}475{,}208\)
signatures through dimension 200, of which 2145 force ordinarity.
The companion verification archive contains the dimension-30 table,
the larger summary, and the replays for both computations.
